
%% ============================================================
%% JOURNAL OF NONPARAMETRIC STATISTICS (Taylor & Francis)
%% Uses interact.cls from Taylor_Francis_Template/
%% SUBMISSION INSTRUCTIONS (Updated 16-02-2026):
%%   - Submit via Taylor & Francis Submission Portal
%%   - Convert this file to PDF for upload ("Manuscript" file)
%%   - Upload LaTeX source files as a single ZIP alongside the PDF
%%   - Peer review is SINGLE-BLIND: include full author details below
%%   - British (-ise) spelling required throughout (already used here)
%%   - Abstract: <= 150 words (unstructured)
%%   - Keywords: 2-5 only
%%   - Include a word count in the manuscript
%% ============================================================
\documentclass{interact}

\usepackage{amsmath,amssymb,amsthm}
\usepackage{mathtools}        % required for \coloneqq
\usepackage{booktabs}         % for \toprule, \midrule, \bottomrule in tables
\usepackage{array}            % extended column types in tables
\usepackage{hyperref}

%% natbib citation style as required by interact.cls / tfcad
\usepackage{natbib}
\bibpunct[, ]{(}{)}{;}{a}{}{,}
\renewcommand\bibfont{\fontsize{10}{12}\selectfont}

%% Theorem-like environments following interact.cls template conventions.
%% All share the [theorem] counter so numbering is sequential within sections.
\theoremstyle{plain}
\newtheorem{theorem}{Theorem}[section]
\newtheorem{lemma}[theorem]{Lemma}
\newtheorem{corollary}[theorem]{Corollary}
\newtheorem{proposition}[theorem]{Proposition}

\theoremstyle{definition}
\newtheorem{definition}[theorem]{Definition}
\newtheorem{assumption}[theorem]{Assumption}

\theoremstyle{remark}
\newtheorem{remark}{Remark}

%% interact.cls: \articletype is optional; use "ARTICLE" for standard papers
\articletype{ARTICLE}

\title{Anti-Modes, Distributional Gaps, and the Fisher Information
Profile: Nonparametric Estimation and Inference}

%% interact.cls author block: use \name{} for names and \affil{} for affiliations.
%% GNST is SINGLE-BLIND: full author details required.
%% Replace the placeholders below with your actual information.
%% Example for one author:
%%   \author{
%%   \name{Jane Smith\thanks{CONTACT Jane Smith.
%%     Email: j.smith@university.ac.uk}}
%%   \affil{Department of Statistics, University X, City, Country}
%%   }
\author{
\name{Samir Orujov%
\thanks{CONTACT Samir Orujov.
Email: sorujov@ada.edu.az}}
\affil{ADA University / ICTA / UNEC, Baku, Azerbaijan;
Universit\'e Bretagne Sud, Vannes, France}
}

\begin{document}
\maketitle

%% ============================================================
\begin{abstract}
We develop a nonparametric theory of distributional gaps centred on
the \emph{Fisher information profile} $I(p)=f(Q(p))^2/[p(1-p)]$,
where $Q$ is the quantile function and $f$ the density.  A
percentile $p^*$ is an \emph{anti-mode percentile} if and only if it
is simultaneously a local minimiser of $f$, a local maximiser of the
quantile density $q(p)=Q'(p)$, and a local maximiser of the asymptotic
quantile estimation variance $\sigma^2_Q(p)$---and, under a mild
first-order condition, a local minimiser of
$I(p)$---providing the first information-theoretic
characterisation of distributional gaps.  We prove $n^{2/5}$
consistency and asymptotic normality of the kernel plug-in estimator
$\hat p^*$, uniform convergence of $\hat I(p)$ at rate
$O_p((nh)^{-1/2}+h^2)$, and bootstrap size control with nontrivial
power against local alternatives shrinking at rate $n^{-2/5}$.
Monte Carlo experiments across twelve data-generating processes
confirm the theoretical rates, and an application to Old Faithful
geyser data illustrates the full inferential output.
\end{abstract}
%% Word count (main text, excluding abstract, references, tables):
%% approximately 2,700 words (texcount after Batch 2 edits; recount before submission).

\begin{keywords}
anti-mode; Fisher information profile; quantile density;
nonparametric testing; bootstrap inference
\end{keywords}

\begin{amscode}
62G07; 62G10; 62G20; 62E20
\end{amscode}

%% ============================================================
\section{Introduction}
\label{sec:intro}

The mode of a probability distribution---the value at which the
density attains a global or local maximum---has occupied a central
position in nonparametric statistics since the foundational work of
\citet{Parzen1962} and \citet{Chernoff1964}.  Mode estimation,
testing for unimodality, and modal clustering have generated a rich
literature spanning kernel methods \citep{Romano1988},
bump-hunting \citep{GoodSilverman1986}, and most recently,
topological data analysis \citep{Chazal2018}.  Foundational work on
kernel-based multimodality testing includes \citet{Silverman1981},
who proposed the critical bandwidth test, and
\citet{MullerSawitzki1991}, who developed excess mass estimates for
distributional gaps; neither provides an estimate of where the
anti-mode lies nor an asymptotic theory for its location.
The anti-mode---the
dual object, a local \emph{minimum} of the density---has received
virtually no systematic theoretical treatment despite its natural
interpretation as a distributional gap or structural fault line.
Existing procedures such as the dip test \citep{HartiganHartigan1985}
and Z-Dip \citep{ZDip2025} test for the \emph{presence} of
multimodality but provide no asymptotic theory for the location of
the gap, no consistent estimator of the anti-mode percentile, and
no confidence interval for its position.

Anti-modes arise in every domain where populations separate into
distinct subgroups with a sparse transition zone.  In income and
wage distributions, sparse income regions impede upward mobility and
correspond to socioeconomic class boundaries
\citep{Orujov2025EcLet}.  In clinical biomarker data,
anti-modes between unimodal sub-populations of diseased and healthy
individuals signal the natural threshold for diagnostic
classification.  In financial return distributions, a sparse zone
separating normal and extreme-loss regimes reflects mixture
structure associated with crisis episodes
\citep{ContTankov2004}.  In all these cases, the location,
intensity, and statistical precision of the anti-mode carry
substantive content that aggregate summary statistics (mean,
variance, Gini coefficient) cannot reveal.

This paper makes four contributions.  \textit{First}, we introduce
the \emph{Fisher information profile} $I(p)$ for a general
univariate distribution and establish equivalent characterisations
of anti-modes via density, quantile density, and Fisher information,
providing the first information-theoretic characterisation of
distributional gaps (Section~\ref{sec:theory}).  \textit{Second},
we derive the asymptotic distribution of the kernel plug-in
anti-mode estimator, showing $n^{2/5}$-consistency and asymptotic
normality, and provide the optimal bandwidth minimising asymptotic
mean squared error (Section~\ref{sec:asymptotics}).
\textit{Third}, we construct a bootstrap test for the null of no
anti-mode and derive its asymptotic size and local power
(Section~\ref{sec:inference}).
\textit{Fourth}, we evaluate the framework across twelve simulation
DGPs spanning Gaussian, heavy-tailed, and multi-modal families
(Section~\ref{sec:simulation}), and illustrate the
full inferential output on the Old Faithful geyser dataset
(Section~\ref{sec:empirical}).  A companion paper
\citep{Orujov2025EcLet} applies this framework to harmonised
income micro-data across 25 countries.

The paper is self-contained and does not require income-specific
objects such as the Lorenz curve; the Lorenz curvature appears as a
corollary of the general theory for non-negative distributions with
finite mean (Remark~\ref{rem:lorenz}).

%% ============================================================
\section{The Fisher Information Profile and Anti-Mode Theory}
\label{sec:theory}

\subsection{Setup and notation}

Let $X$ be a real-valued random variable with distribution function
$F$, quantile function $Q(p)=F^{-1}(p)=\inf\{x:F(x)\ge p\}$ for
$p\in(0,1)$, and Lebesgue density $f=F'$.  We impose the following
standing conditions.

\begin{assumption}[Regularity]\label{ass:regularity}
\begin{enumerate}
  \item[(A1)] $f$ is continuous and strictly positive on the support
    of $X$, and $Q$ is twice continuously differentiable on
    $(0,1)$.
  \item[(A2)] The quantile density $q(p)\coloneqq Q'(p)=
    1/f(Q(p))$ is bounded away from zero and infinity on
    $[\delta,1-\delta]$ for every $\delta\in(0,\tfrac{1}{2})$.
  \item[(A3)] $q$ is three times continuously differentiable on
    $(\delta,1-\delta)$.
\end{enumerate}
\end{assumption}

These conditions are standard in the quantile-process literature
\citep{Serfling1980, vanderVaartWellner1996} and are satisfied
by all commonly used parametric families (normal, $t$, logistic,
exponential, Pareto, Weibull) on their supports (mixtures of such
families may violate A3 at the anti-mode when the mixture valley is
shallow and asymmetric; DGP~B3 below illustrates this).

\subsection{The Fisher information profile}

The asymptotic distribution of the sample quantile is classical:
$\sqrt{n}(\hat Q_n(p)-Q(p))\xrightarrow{d}
N(0,\,p(1-p)/f(Q(p))^2)$ \citep{Serfling1980}.  Define the
\emph{asymptotic quantile variance profile}
\begin{equation}
\sigma^2_Q(p) \;=\; \frac{p(1-p)}{f(Q(p))^2}
               \;=\; p(1-p)\,q(p)^2,
\label{eq:qvar}
\end{equation}
and the \emph{Fisher information profile}
\begin{equation}
I(p) \;=\; \frac{f(Q(p))^2}{p(1-p)}
      \;=\; \frac{1}{p(1-p)\,q(p)^2}
      \;=\; \frac{1}{\sigma^2_Q(p)}.
\label{eq:infoprofile}
\end{equation}
Equation~\eqref{eq:infoprofile} has a precise statistical
interpretation: $I(p)^{-1}$ is the asymptotic Cram\'{e}r--Rao lower
bound for unbiased estimation of $Q(p)$ from $n$ i.i.d.\
observations \citep{LehmannCasella1998}.  It also coincides with
the Fisher information for the one-parameter location family
$F(x-\theta_p)$ evaluated at $\theta_p=Q(p)$
\citep{vanderVaartWellner1996}.

The profile $p\mapsto I(p)$ is a functional descriptor of the
entire shape of $F$: it is large where $f(Q(p))$ is large (dense
regions, tight quantile spacing) and small where $f(Q(p))$ is small
(sparse regions, wide quantile spacing).

\subsection{Anti-modes: definitions and equivalent characterisations}

\begin{definition}[Anti-mode percentile]\label{def:antimode}
A percentile $p^*\in(\delta,1-\delta)$ is an
\emph{anti-mode percentile} of $F$ if any one---hence all---of the
following equivalent conditions holds:
\begin{enumerate}
  \item[(i)] $x^*=Q(p^*)$ is a strict local minimiser of $f$
    \textnormal{(density anti-mode)};
  \item[(ii)] $p^*$ is a strict local maximiser of the quantile
    density $q(p)$;
  \item[(iii)] $p^*$ is a strict local maximiser of
    $\sigma^2_Q(p)/[p(1-p)] = q(p)^2$.
\end{enumerate}
The \emph{anti-mode intensity} $\eta^*$ is the prominence of the
local maximum of $q$ at $p^*$, defined as
$\eta^*=q(p^*)-\max(q_L,q_R)$ where $q_L$ and $q_R$ are the
nearest local minima of $q$ to the left and right of $p^*$.
\end{definition}

The equivalence of conditions (i)--(iii) is immediate from the
identity $q(p)=1/f(Q(p))$ and \eqref{eq:infoprofile}.

\begin{theorem}[Equivalent characterisations of anti-modes]\label{thm:equivalence}
Under Assumption~\ref{ass:regularity}, conditions \textnormal{(i)},
\textnormal{(ii)}, and \textnormal{(iii)} of
Definition~\ref{def:antimode} are mutually equivalent.
\end{theorem}

\begin{proof}
(i)$\Leftrightarrow$(ii): $x^*=Q(p^*)$ is a strict local minimiser
of $f$ if and only if $1/f(Q(p^*))=q(p^*)$ is a strict local
maximiser of $q$, since $t\mapsto 1/t$ is strictly decreasing on
$(0,\infty)$.

(ii)$\Leftrightarrow$(iii): Since $q>0$, the map $t\mapsto t^2$ is
strictly increasing on $(0,\infty)$, so strict local maxima of $q$
coincide exactly with strict local maxima of $q^2=\sigma^2_Q(p)/[p(1-p)]$.
$\square$
\end{proof}

\begin{corollary}[Fisher information characterisation]\label{cor:fisher}
Under Assumption~\ref{ass:regularity}, $p^*$ is a strict local
minimiser of $I(p)=1/[p(1-p)q(p)^2]$ if and only if $p^*$ is a
strict local maximiser of $\sigma^2_Q(p)=p(1-p)q(p)^2$.
This Fisher information characterisation coincides with conditions
\textnormal{(i)}--\textnormal{(iii)} of Definition~\ref{def:antimode}
whenever the first-order condition
\begin{equation}\label{eq:foc}
  \frac{d}{dp}\bigl[p(1-p)q(p)^2\bigr]\Big|_{p^*}
  = (1-2p^*)q(p^*)^2 + 2p^*(1-p^*)q(p^*)q'(p^*) = 0
\end{equation}
holds at $p^*$.  In particular, conditions \textnormal{(i)}--\textnormal{(iii)}
and the Fisher information characterisation are jointly satisfied at
every anti-mode of a distribution symmetric about its median (where
$p^*=\tfrac{1}{2}$), and approximately so for near-symmetric
bimodal distributions.
\end{corollary}

\begin{proof}
The equivalence $I(p)$ local min $\Leftrightarrow$ $\sigma^2_Q(p)$ local max
follows from $I=1/\sigma^2_Q$ and strict monotonicity of $t\mapsto 1/t$.
Differentiating $I$ at $p^*$ with $q'(p^*)=0$ gives
$I'(p^*)=-(1-2p^*)q(p^*)^2/[p^*(1-p^*)q(p^*)^2]^2$, which equals zero
if and only if \eqref{eq:foc} holds at $p^*$.
The second-derivative condition then confirms a strict local minimum of $I$.
$\square$
\end{proof}

\begin{remark}[Lorenz corollary]\label{rem:lorenz}
For non-negative $X$ with finite mean $\mu$, the Lorenz curve
satisfies $L''(p) = 1/[\mu f(Q(p))] = q(p)/\mu$. Therefore all
local maxima of $L''(p)$ coincide with local maxima of $q(p)$ and
hence with anti-modes. Income-distributional glass barriers are a special case of
Definition~\ref{def:antimode}; a companion paper \citep{Orujov2025EcLet}
applies this framework to income distributions across countries and
establishes the glass-barrier concept empirically, while the present paper
provides the statistical foundations.
\end{remark}

\subsection{Axiomatic properties}

\begin{proposition}[Scale and location invariance]
\label{prop:scaleloc}
Let $\tilde X = aX+b$ for $a>0,b\in\mathbb{R}$.  Then:
\begin{enumerate}
  \item[(i)] The set of anti-mode percentiles of $\tilde X$ equals
    that of $X$.
  \item[(ii)] Anti-mode intensities are unchanged by location
    shifts ($a=1$, arbitrary $b$) and scale as $a$ under
    scale changes ($b=0$, arbitrary $a$).
  \item[(iii)] The Fisher information profile satisfies
    $\tilde I(p)=I(p)/a^2$.
\end{enumerate}
\end{proposition}

\begin{proof}
$\tilde Q(p)=aQ(p)+b$ and $\tilde f(\tilde x)=f((\tilde
x-b)/a)/a$, so $\tilde q(p)=\tilde
Q'(p)=aQ'(p)=aq(p)$. Local maxima of $\tilde q=aq$ occur at the
same percentiles as those of $q$; prominence scales by $a$.
Substituting into \eqref{eq:infoprofile} gives
$\tilde I(p)=\tilde f(\tilde Q(p))^2/[p(1-p)]=
[f(Q(p))/a]^2/[p(1-p)]=I(p)/a^2$. $\square$
\end{proof}

\begin{proposition}[Mixture representation]
\label{prop:mixture}
Let $F=\pi_1 F_1+(1-\pi_1)F_2$ be a two-component mixture with
weights $\pi_1\in(0,1)$.  If $F_1$ and $F_2$ are unimodal and
their supports overlap, then $F$ generically has an anti-mode at a
percentile $p^*$ in the overlap region; the intensity $\eta^*$
increases in the separation between the component modes and
decreases in $\pi_1(1-\pi_1)$ (i.e., more balanced mixtures have
stronger barriers).
\end{proposition}

\begin{proof}
The mixture density is $f=\pi_1 f_1 + (1-\pi_1)f_2$.  In the
overlap region between the modes $m_1$ and $m_2$ (with $m_1<m_2$),
$f_1$ is decreasing and $f_2$ is increasing.  Under standard
smoothness and non-degeneracy conditions, $f$ attains a local
minimum at some $x^*\in(m_1,m_2)$.  The depth of this minimum
(and hence the intensity $\eta^*$) increases with $|m_1-m_2|$
and with $\min(\pi_1,1-\pi_1)$. $\square$
\end{proof}

%% ============================================================
\section{Asymptotic Theory}
\label{sec:asymptotics}

\subsection{Estimation setup}

Given i.i.d.\ observations $X_1,\ldots,X_n$ from $F$, define:
\begin{itemize}
  \item the empirical quantile $\hat Q_n(p) = X_{(\lceil np\rceil)}$;
  \item the kernel density estimator $\hat f_h(x) =
    (nh)^{-1}\sum_{i=1}^n K((x-X_i)/h)$ with symmetric kernel $K$
    satisfying $\int K=1$, $\int u^2 K(u)\,du=\kappa_2<\infty$,
    and bandwidth $h=h_n\to 0$, $nh\to\infty$;
  \item the plug-in quantile density estimator
    $\hat q(p) = 1/\hat f_h(\hat Q_n(p))$;
  \item the estimated Fisher information profile
    $\hat I(p)=1/[p(1-p)\hat q(p)^2]$.
\end{itemize}
Anti-mode percentiles are estimated as
$\hat p^* = \operatorname{argmax}_{p\in[\delta,1-\delta]}\hat q(p)$
restricted to isolated local maxima with prominence exceeding a
threshold $\tau>0$.

\subsection{Uniform convergence of $\hat I(p)$}

\begin{theorem}[Uniform convergence]\label{thm:uniform}
Under Assumption~\ref{ass:regularity} and $h\to 0$, $nh\to\infty$:
\begin{equation}
\sup_{p\in[\delta,1-\delta]}\bigl|\hat I(p)-I(p)\bigr|
\;=\; O_p\!\bigl((nh)^{-1/2}+h^2\bigr).
\end{equation}
\end{theorem}

\begin{proof}[Proof sketch]
By a Taylor expansion of $\hat I$ around $I$,
$|\hat I(p)-I(p)|\lesssim |(\hat q(p)/q(p))^2-1|
\lesssim 2|\hat q(p)/q(p)-1|+(\hat q(p)/q(p)-1)^2$.
Since $\hat q(p)=1/\hat f_h(\hat Q_n(p))$ and $q(p)=1/f(Q(p))$,
the key step is to control
$\sup_{p}|\hat f_h(\hat Q_n(p))-f(Q(p))|$.
Writing $\hat f_h(\hat Q_n(p))-f(Q(p))
= [\hat f_h(\hat Q_n(p))-\hat f_h(Q(p))]
  + [\hat f_h(Q(p))-f(Q(p))]$,
the second bracketed term is controlled uniformly by
$O_p((nh)^{-1/2}+h^2)$ using standard KDE uniform convergence
results \citep{Parzen1962, Jiang2017}.  The first term is bounded
by $\sup_x|\hat f_h'(x)|\cdot|\hat Q_n(p)-Q(p)|
=O_p((nh^3)^{-1/2})\cdot O_p(n^{-1/2})=o_p(1)$ under the
stated bandwidth conditions.  Since $q>c>0$ on
$[\delta,1-\delta]$ by A2, the result follows from the delta
method applied to $q\mapsto 1/q^2$. $\square$
\end{proof}

\subsection{Convergence rate and asymptotic normality of $\hat p^*$}

Let $p^*$ be an isolated strict local maximiser of $q$ with
$q''(p^*)<0$.  Define $\gamma^2 = -q''(p^*)>0$.

\begin{theorem}[Asymptotic normality of anti-mode estimator]
\label{thm:asymptotics}
Under Assumption~\ref{ass:regularity}, $h\sim c_0 n^{-1/5}$ for
some $c_0>0$, and $q'''$ continuous in a neighbourhood of $p^*$:
\begin{equation}
n^{2/5}\bigl(\hat p^* - p^* - B_n\bigr)
\;\xrightarrow{d}\; N\!\left(0,\,V\right),
\label{eq:asymnorm}
\end{equation}
where the asymptotic bias is
\begin{equation}
B_n \;=\; -\frac{\kappa_2\,c_0^2\,q'''(p^*)}{2\gamma^2\,f(Q(p^*))^3}
           \cdot n^{-2/5}\,(1+o(1)),
\label{eq:bias}
\end{equation}
with $q'''(p^*) = (d^3q/dp^3)|_{p=p^*}$ the third derivative of the quantile
density at $p^*$, and the asymptotic variance is
\begin{equation}
V \;=\; \frac{\|K'\|_2^2\,q(p^*)^5}{c_0^3\,\gamma^4},
\label{eq:variance}
\end{equation}
where $\|K'\|_2^2 = \int [K'(u)]^2\,du$ is the roughness of the kernel
derivative.  Note that $V$ is a fixed constant (independent of $n$).
\end{theorem}

\begin{remark}[Practical inference]\label{rem:practical}
Because $B_n = O(n^{-2/5})$, two strategies eliminate the bias from confidence
interval construction: (a) \emph{undersmoothing} with $h = o(n^{-1/5})$ so
that $n^{2/5}B_n\to 0$ and the standard $n^{2/5}(\hat p^* - p^*)
\xrightarrow{d}N(0,V)$ holds without bias correction; or (b) a
\emph{plug-in bias estimate} obtained by evaluating $\hat q'''(p^*)$ from a
third-order kernel derivative estimator at bandwidth $h$.  In our simulations
we use strategy (a) with $h = 0.636\hat\sigma n^{-1/5}$
(undersmoothed Silverman rule, equal to $0.6\times 1.06\,\hat\sigma\,n^{-1/5}$).
\end{remark}

\begin{proof}[Proof sketch]
The estimator $\hat p^*$ is the argmax of $\hat q(p)$ near $p^*$.
A Taylor expansion gives $\hat q(p)\approx \hat q(p^*)
+\hat q'(p^*)(p-p^*)+\frac{1}{2}\hat q''(p^*)(p-p^*)^2$.  Setting
the derivative to zero and solving yields
$\hat p^*-p^*\approx -\hat q'(p^*)/\hat q''(p^*)$.
The numerator $\hat q'(p^*)$ has bias $B_n\cdot\gamma^2$ and variance
of order $(nh^3)^{-1}$ (derivative KDE variance), while $\hat
q''(p^*)\xrightarrow{p} q''(p^*)=-\gamma^2$.  Rescaling by
$n^{2/5}$ (the rate equating bias $\sim h^2\sim n^{-2/5}$ and
standard deviation $\sim(nh^3)^{-1/2}\sim n^{-2/5}$ at $h\sim
n^{-1/5}$) and applying Slutsky's theorem gives
\eqref{eq:asymnorm}--\eqref{eq:variance}.  Full details are in
Appendix~\ref{app:proofs}. $\square$
\end{proof}

\begin{corollary}[Optimal bandwidth]\label{cor:optbw}
The MSE-optimal bandwidth for anti-mode percentile estimation is
\begin{equation}
h^*_{\mathrm{AMSE}} \;=\;
\left(\frac{V_0}{2\,B_0^2}\right)^{1/5} n^{-1/5},
\end{equation}
where $V_0$ and $B_0$ are the leading constants in
\eqref{eq:bias}--\eqref{eq:variance} evaluated at $p^*$. Under
$h=h^*_{\mathrm{AMSE}}$, $\mathrm{AMSE}(\hat p^*)=O(n^{-4/5})$.
\end{corollary}

\begin{remark}
The $n^{-4/5}$ AMSE rate matches the minimax-optimal rate for
estimating the location of a mode (or anti-mode) of a smooth
density \citep{Romano1988}, confirming that the plug-in estimator
is rate-optimal.
\end{remark}

%% ============================================================
\section{Inference: Detection, Confidence Sets, and Testing}
\label{sec:inference}

\subsection{Anti-mode detection algorithm}

Given $\hat q(p)$ on a grid $\{p_j\}_{j=1}^M$ over
$[\delta,1-\delta]$:
\begin{enumerate}
  \item Locate all interior local maxima of $\hat q$ (candidates).
  \item Retain candidates whose prominence exceeds $\tau$ times
    the flanking valley height.
    Specifically, let $L_j=\min_{i<j}\hat q(p_i)$ and
    $R_j=\min_{i>j}\hat q(p_i)$, and define the valley height
    $V_j=\max(L_j,R_j)$.  A local maximum at $p_j$ is retained if
    $\hat q(p_j)-V_j\geq \tau\,V_j$, i.e.\ the peak exceeds
    $V_j$ by at least a fraction~$\tau$.
    This local-prominence criterion is scale-invariant and
    converges to a finite population value as $n\to\infty$,
    since $\hat q(p)\to q(p)$ uniformly on $[\delta,1-\delta]$.
  \item Merge candidates within $\kappa$ percentile points,
    retaining the higher peak.
  \item Report surviving candidates $\{\hat p^*_j\}$ with
    quantile levels $\hat x^*_j=\hat Q(\hat p^*_j)$ and
    intensities $\hat\eta^*_j$.
\end{enumerate}

\subsection{Bootstrap confidence sets}

Because $\hat p^*$ is an argmax functional, the delta method does
not apply.  We use the nonparametric bootstrap with $B$ resamples
to construct equal-tailed $95\%$ confidence intervals.  When a
bootstrap resample detects a different number of anti-modes than
the original estimate, we use a matching algorithm based on
Hausdorff distance on the set of detected percentiles.

\begin{theorem}[Bootstrap consistency]\label{thm:bootstrap}
Under Assumption~\ref{ass:regularity} and the conditions of
Theorem~\ref{thm:asymptotics}, the bootstrap distribution of
$n^{2/5}(\hat p^{*,b}-\hat p^*)$ consistently approximates the
distribution of $n^{2/5}(\hat p^*-p^*)$ conditional on
$X_1,\ldots,X_n$, in the bounded Lipschitz metric.
\end{theorem}

\begin{proof}
The result follows from bootstrap consistency for smooth
functionals of the empirical process \citep{vanderVaartWellner1996}
combined with the delta method argument in the proof of
Theorem~\ref{thm:asymptotics}.  Details are in
Appendix~\ref{app:proofs}. $\square$
\end{proof}

\subsection{Test for no anti-mode}

Consider testing $H_0: f \text{ is unimodal on }(Q(\delta),
Q(1-\delta))$ against $H_1: f$ has at least one anti-mode.  Define
the test statistic
\begin{equation}
T_n = \max_{p\in[\delta,1-\delta]}
  \bigl[\hat q(p) - \hat q_0(p)\bigr]_+,
\end{equation}
where $\hat q_0$ is the quantile density under a parametric
unimodal null (e.g., fitted log-normal).  Large $T_n$ indicates
evidence against $H_0$.

\begin{theorem}[Size control]\label{thm:test}
Under $H_0$ and the conditions of Theorem~\ref{thm:uniform}, the
parametric bootstrap approximation to the null distribution of $T_n$
achieves asymptotic level $\alpha$ for any $\alpha\in(0,1)$.
Under local alternatives with anti-mode intensity
$\eta_n\sim c\,n^{-2/5}$ for $c>c_{\min}(\alpha)$, the power of the
test converges to one.
\end{theorem}

\begin{proof}[Proof sketch]
\textit{Size control.}  Under $H_0$, $\hat q_0$ is fitted to a
correctly specified unimodal family, so $\sup_p|\hat q_0(p)-q_0(p)|
=O_p(n^{-1/2})$ by standard parametric consistency.  By
Theorem~\ref{thm:uniform}, $\sup_p|\hat q(p)-q(p)|
=O_p((nh)^{-1/2}+h^2)$.  Under $H_0$, $q=q_0$ so
$T_n=\sup_p[\hat q(p)-\hat q_0(p)]_+ \xrightarrow{p} 0$.
The parametric bootstrap resamples from $\hat F_0$ (the fitted null
model) and recomputes $T_n^*$ on each resample.  Since both the
statistic and its bootstrap analogue converge uniformly in $p$,
the bootstrap distribution of $T_n^*$ consistently approximates the
null distribution of $T_n$ in total variation; see
\citet{vanderVaartWellner1996}, Theorem 23.9.

\textit{Power.}  Under a local alternative with $\eta_n=q(p^*)-
\max(q_L,q_R)\sim cn^{-2/5}$, by Theorem~\ref{thm:asymptotics}
the estimator $\hat q(p^*)$ detects the local maximum with
probability tending to one for $c>c_{\min}(\alpha)$, where
$c_{\min}(\alpha)$ depends on $\alpha$, $V^{1/2}$, and the critical
value of the bootstrap null distribution.  Hence
$T_n\xrightarrow{p}\infty$ relative to the bootstrap critical value,
giving power tending to one. $\square$
\end{proof}

%% ============================================================
\section{Monte Carlo Study}
\label{sec:simulation}

\subsection{Design}

We evaluate performance across twelve DGPs falling into three
families:

\noindent\textbf{Family A -- Gaussian mixtures (6 DGPs):}
$F=\pi N(\mu_1,\sigma_1^2)+(1-\pi)N(\mu_2,\sigma_2^2)$ with
varying separation $\Delta=|\mu_2-\mu_1|/\bar\sigma$ and mixing
weight $\pi$.  DGP A1: $\Delta=2.6$, $\pi=0.5$ (weak barrier);
DGP A2: $\Delta=2.8$, $\pi=0.5$ (moderate); DGP A3: $\Delta=3$,
$\pi=0.5$ (sharp); DGP A4: $\Delta=3$, $\pi=0.3$ (unequal
weights); DGP A5: $\Delta\approx2.7$, $\pi=0.5$, $\sigma_2=2\sigma_1$
(unequal variances); DGP A6: unimodal $N(0,1)$ (size baseline).
All DGPs A1--A5 have a genuine density anti-mode; the bimodality
threshold for equal-variance Gaussian mixtures is $\Delta=2$, so
all five are strictly above it (or, for A5, above the corresponding
threshold for unequal variances).

\noindent\textbf{Family B -- Skew and heavy-tailed mixtures (3 DGPs):}
DGP B1: mixture of skew-normal distributions; DGP B2: mixture of
Student-$t_3$ components (heavy tails); DGP B3: mixture of
$\mathrm{Weibull}(\alpha_1,\lambda)$ and
$\mathrm{Weibull}(\alpha_2,\lambda)$ (asymmetric, non-negative support,
relevant for income and survival data).

\noindent\textbf{Family C -- Multi-modal (3 DGPs):}
DGP C1: three-component Gaussian mixture with two anti-modes;
DGP C2: four-component mixture with three anti-modes; DGP C3:
``bathtub'' shape (U-shaped density, single deep anti-mode at
the median).

All experiments use $n\in\{250,500,1000,5000\}$,
$R=500$ replications, and $B=299$ bootstrap draws per
replication.

\subsection{Results}

Table~\ref{tab:bias_rmse} reports bias and RMSE of $\hat p^*$ (in
percentile units) for the primary anti-mode.  Bias is small and
shrinks with $n$ for Gaussian and heavy-tailed mixtures (Families~A
and~B2), consistently with Theorem~\ref{thm:asymptotics}.  DGP~B3
(Weibull mixture) is a notable exception: its anti-mode lies in a
shallow, asymmetric valley, and the kernel plug-in exhibits a
persistent positive bias that grows with $n$, suggesting the
smoothness conditions of Theorem~\ref{thm:asymptotics} are
effectively violated for this design.
RMSE declines at approximately $n^{-2/5}$: regressing
$\log(\mathrm{RMSE})$ on $\log(n)$ yields slopes whose 95\%
bootstrap confidence intervals include $-0.40$ for DGPs~A2, A3,
and B2 (Table~\ref{tab:rates}).  The multi-modal designs C1 and
C3 converge \emph{faster} than $n^{-2/5}$ (slopes around $-0.49$),
consistent with their sharply separated modes.  DGP~B3 converges
at only $-0.11$, reflecting the bias noted above.

Table~\ref{tab:coverage} reports bootstrap coverage rates and
detection rates.  Coverage \emph{exceeds} the nominal 95\% level at
$n=250$--1000 for most bimodal DGPs (conservative intervals, consistent
with the percentile bootstrap overshooting for smooth argmax functionals
\citep{HallYork2001, Efron1987}), and settles near 90--94\% at
$n=5000$ where estimator bias begins to dominate interval width.
Detection rates
generally increase with $n$ for well-separated mixtures: DGP~A3
($\Delta=3$) exceeds 93\% at $n=5000$ and DGP~B2 (heavy-tailed)
reaches 97\%.  For the weak barrier (DGP~A1, $\Delta=2.6$),
detection fluctuates near 70\% and dips to 62\% at $n=5000$,
as the improved KDE resolution reveals that the valley barely
exceeds the prominence threshold~$\tau$.
This non-monotone pattern is a feature of the fixed-$\tau$ design:
as $n\to\infty$ the KDE converges uniformly to the true density,
and for barriers with intensity close to the threshold
the estimator correctly classifies the distribution as
effectively unimodal at that resolution level.
A data-driven, asymptotically vanishing choice of $\tau$ is
a direction for future work.
Under the unimodal null (DGP~A6), the detection algorithm's false
discovery rate is 13.0\% at $n=500$ and 9.6\% at $n=1000$, declining
to 1.0\% at $n=5000$; the elevated small-$n$ rate reflects the fixed
prominence threshold $\tau$ admitting spurious KDE peaks that vanish
as the density estimate stabilises.
DGP~B3 (Weibull mixture) exhibits declining detection
(63\%~$\to$~23\% across $n=250$ to $n=5000$) and undercoverage at large~$n$; both are
consequences of the shallow, asymmetric valley noted above, and
results for this design should be interpreted with caution.

Table~\ref{tab:power} compares the anti-mode test to Hartigan's dip
test \citep{HartiganHartigan1985} and the Z-Dip \citep{ZDip2025}
for DGP~A2 ($\Delta=2.8\sigma$)---a genuinely difficult case:
\citet{ZDip2025} validate their statistic exclusively on mixtures
with component separation $\geq 8\sigma$, where all tests achieve
perfect power and meaningful discrimination is impossible.
The three tests address complementary questions.  The dip test and
Z-Dip are designed for binary multimodality detection and optimised
accordingly; neither provides an estimate of \emph{where} the
anti-mode lies.  The anti-mode test is a \emph{location} test
that, upon rejection, returns a consistent estimator $\hat p^*$
with a calibrated 95\% confidence interval and intensity measure
$\hat\eta^*$---output the other two tests cannot supply.  This
complementarity has a cost: the sup-norm test statistic
$T_n=\sup_p[\hat q(p)-\hat q_0(p)]_+$ is conservative at
moderate sample sizes, yielding near-zero rejection rates below
$n=5000$ for DGP~A2 before reaching 68.8\% at $n=5000$, well
below the dip test (99.8\%).  The conservatism arises because the
sup-norm statistic must exceed the bootstrap critical value
simultaneously across all percentiles $p\in[\delta,1-\delta]$; at
moderate $n$, KDE variance over the full grid inflates both the
observed $T_n$ and the null distribution, masking the localised
signal at $\hat p^*$.  A pointwise variant of the test, evaluated
only at $\hat p^*$ rather than over the full grid, would sacrifice
the omnibus property but substantially improve power against
specific alternatives; this is a direction for future work.
Practitioners requiring only a
detection decision should prefer the dip test; the anti-mode test
is intended for inference on the \emph{location} of the gap.

%% ============================================================
\section{Empirical Illustration: Old Faithful Eruptions}
\label{sec:empirical}

As a compact demonstration, we apply the framework to the Old
Faithful geyser data \citep{AzzaliniBowman1990}---272 eruption
durations (minutes)---a canonical bimodal dataset in nonparametric
statistics.  The detection algorithm identifies a single anti-mode at
$\hat p^*=0.357$, corresponding to $\hat x^*=Q(\hat p^*)=3.00$
minutes; this lies squarely in the well-known gap between short
(${\sim}2$~min) and long (${\sim}4.5$~min) eruption types.  The
bootstrap 95\% confidence interval is $[0.299,\,0.414]$ for $p^*$,
or equivalently $[2.30,\,3.72]$ minutes in the original scale.
Both the sup-norm test ($T_n$) and the dip test reject the unimodal
null at the 5\% level ($p_{\text{dip}}<0.0001$).  That $T_n$ rejects
at $n=272$ here while showing zero power at $n=250$ for DGP~A2
reflects the much sharper Old Faithful gap: the eruption-type
separation is consistent with DGP~A3 or~C3, where detection rates
exceed 93\% even at small~$n$.  The Fisher
information profile at the anti-mode is $\hat I(\hat p^*)=0.007$,
reflecting high estimation uncertainty at the sparse transition zone.
This example illustrates the full inferential output of the
framework---detection, location, confidence set, and
information-theoretic characterisation---on a single real dataset
with $n=272$.

%% ============================================================
\section{Discussion}
\label{sec:discussion}

We have established a complete theory for the nonparametric
estimation and inference of distributional gaps through the Fisher
information profile.

\textit{Comparison with detection-only procedures.}
The dip test \citep{HartiganHartigan1985} and its recent
standardised extension Z-Dip \citep{ZDip2025} answer the binary
question ``is there a gap?''  They provide no asymptotic theory
for the location of the gap, no consistent estimator, and no
confidence interval.  The Z-Dip, in particular, is validated
exclusively on mixtures with component separation $\geq 8\sigma$
\citep{ZDip2025}, where every reasonable test achieves perfect
power; the genuinely difficult borderline regime (such as DGP~A2
at $\Delta=2.8\sigma$) is absent from their study.  The present
framework answers ``where is the gap, how wide is it, and how
precisely can we locate it?''---a strictly richer inferential
output, obtained at the cost of some power relative to the
dedicated detection procedures.

Several further directions for future work are worth
highlighting.

\textit{Multivariate extension.}  The univariate quantile density
does not extend to dimension $d>1$ without ordering; however,
multivariate anti-modes can be defined as local minima of the
density $f$ and characterised through the Hessian of
$\log f$ at $x^*$ (negative definite at a mode, positive
definite at an anti-mode).  The connection to Fisher information
in this setting involves the score function
$\nabla\log f(X-\theta)$ and its covariance matrix, which
generalises $I(p)$ to a matrix-valued profile.

\textit{Time-varying distributions.}  For time series where
$F_t$ evolves over time, anti-mode percentiles become
time-indexed: $p^*(t)=\operatorname{argmax}_p q_t(p)$.
Non-stationary KDE methods \citep{FanGijbels1996} can be
adapted to track the evolution of $p^*(t)$, providing a
diagnostic for structural change in distributional shape.

\textit{Optimal testing.}  Theorem~\ref{thm:test} establishes
asymptotic size control and local power; the question of
whether the proposed test is asymptotically optimal (in the
minimax sense) against specific classes of local alternatives
is left for future work.

\textit{Connection to topological methods.}  Persistent homology
provides a topological analogue of anti-mode intensity via
persistence diagrams \citep{Chazal2018}: the 0-dimensional
persistence of a density anti-mode corresponds exactly to our
intensity $\eta^*$.  Making this correspondence rigorous and
exploiting the stability theory of persistent homology for
inference is a promising direction.

%% ============================================================
\section*{Acknowledgements}
The author thanks the editor and anonymous referees for their time.

%% ============================================================
\section*{Disclosure statement}
The author reports no competing interests to declare.

%% ============================================================
%% GNST REQUIRED: Funding details.
%% Replace with your actual funding information, or state:
%%   "This research received no specific grant from any funding
%%    agency in the public, commercial, or not-for-profit sectors."
\section*{Funding}
This research received no specific grant from any funding agency in the
public, commercial, or not-for-profit sectors.

%% ============================================================
\section*{ORCID}
Samir Orujov \url{https://orcid.org/0009-0004-9708-2109}

%% ============================================================
\section*{CRediT author statement}
\textbf{Samir Orujov}: Conceptualization, Methodology, Software,
Formal analysis, Investigation, Writing -- Original Draft,
Writing -- Review \& Editing, Visualization.

%% ============================================================
\section*{Data availability statement}
%% GNST REQUIRED: Describe where data supporting the results can be found.
Replication code for all Monte Carlo experiments reported in this paper
is available in an online repository at
\url{https://github.com/sorujov/Antimode-JNPS}.
The Old Faithful geyser data used in Section~\ref{sec:empirical} are
publicly available in the R \texttt{datasets} package
\citep{AzzaliniBowman1990}.

%% ============================================================
\appendix
\section{Proofs}
\label{app:proofs}

\subsection*{Proof of Theorem~\ref{thm:asymptotics} (full detail)}

Let $g(p)=q(p)=1/f(Q(p))$.  By the second-order condition,
$g''(p^*)=-\gamma^2<0$.  The estimator satisfies
$\hat g'(\hat p^*)=0$, so a Taylor expansion around $p^*$ gives
\[
0 = \hat g'(p^*) + \hat g''(p^*)(\hat p^*-p^*) + O_p((\hat p^*-p^*)^2).
\]
Hence
\[
\hat p^*-p^* = -\frac{\hat g'(p^*)}{\hat g''(p^*)}(1+o_p(1)).
\]
\textit{Step 1: Denominator.}
$\hat g''(p^*)\xrightarrow{p}g''(p^*)=-\gamma^2$ by uniform
convergence of $\hat q''$ to $q''$ (Theorem~\ref{thm:uniform}
applied to the second derivative of $\hat q$).

\textit{Step 2: Numerator --- bias.}
Set $x^*=Q(p^*)$.  By the chain rule,
\[
\hat g'(p) \;=\; \frac{d}{dp}\hat q(p)
 \;=\; -\frac{\hat f_h'(\hat Q_n(p))}{\hat f_h(\hat Q_n(p))^2}
       \cdot\frac{d\hat Q_n}{dp}(p).
\]
To leading order, $d\hat Q_n/dp \approx \hat q(p)=1/\hat f_h(\hat Q_n(p))$,
so
\[
\hat g'(p^*) \;\approx\; -\frac{\hat f_h'(x^*)}{\hat f_h(x^*)^3}.
\]
The bias of $\hat f_h'(x^*)$ is $\tfrac{1}{2}\kappa_2 h^2 f'''(x^*)$
\citep{Parzen1962}.  We convert $f'''(x^*)$ to $q$-derivatives as
follows.  From $f(Q(p))=1/q(p)$, differentiating with respect to $p$
and using $Q'(p)=q(p)$ gives $f'(Q(p))=-q'(p)/q(p)^3$.  A second
differentiation gives $f''(Q(p))=[-q''(p)q(p)^3+3q'(p)^2q(p)^2]/q(p)^6
=[-q''(p)+3q'(p)^2/q(p)]/q(p)^3$.  A third differentiation, evaluated
at $p^*$ where $q'(p^*)=0$ (so all terms involving $q'$ vanish), yields
\[
f'''(Q(p^*))\cdot q(p^*) = -q'''(p^*)/q(p^*)^3,
\qquad\text{i.e.,}\quad
f'''(x^*) = -q'''(p^*)/q(p^*)^4.
\]
Applying the delta method to the map $y\mapsto -y/f(x^*)^3$
gives
\[
\operatorname{Bias}(\hat g'(p^*))
\;\approx\; -\frac{\kappa_2 h^2 f'''(x^*)}{2f(x^*)^3}
\;=\; -\frac{\kappa_2 c_0^2 q'''(p^*)}{2f(x^*)^3}\,n^{-2/5},
\]
consistent with \eqref{eq:bias}.

\textit{Step 3: Numerator --- variance.}
Since $q'(p^*)=0$ at the anti-mode, the stochastic fluctuation of
$\hat g'(p^*)$ comes entirely from the randomness of $\hat f_h'(x^*)$.
The variance of the KDE derivative is \citep{Romano1988}
\[
\operatorname{Var}(\hat f_h'(x^*)) \;=\; \frac{\|K'\|_2^2\,f(x^*)}{nh^3}.
\]
Propagating through the delta method for the map $y\mapsto-y/f(x^*)^3$:
\[
\operatorname{Var}(\hat g'(p^*))
\;\approx\; \frac{\|K'\|_2^2\,f(x^*)}{nh^3\,f(x^*)^6}
\;=\; \frac{\|K'\|_2^2\,q(p^*)^5}{nh^3}.
\]
With $h=c_0 n^{-1/5}$ we get $nh^3 = c_0^3 n^{2/5}$, so
$\operatorname{Var}(n^{2/5}\hat g'(p^*))
\to \|K'\|_2^2 q(p^*)^5/c_0^3$,
a finite constant.  Dividing by $\gamma^4$ (from the denominator)
yields the asymptotic variance $V$ in \eqref{eq:variance}.  The
Lindeberg--Feller CLT applies since $\hat f_h'(x^*)$ is an
$L^2$-bounded triangular array, completing the proof. $\square$

\subsection*{Proof of Theorem~\ref{thm:bootstrap}}

The KDE $\hat f_h$ is a smooth functional of the empirical measure
$\hat{\mathbb{P}}_n$ (it is the convolution $\hat{\mathbb{P}}_n * K_h$,
linear in $\hat{\mathbb{P}}_n$), and the plug-in quantile density
$\hat q(p)=1/\hat f_h(\hat Q_n(p))$ is a Hadamard differentiable
map of $\hat{\mathbb{P}}_n$ at every $p$ where $f(Q(p))>0$
(Assumption A2), by the chain rule for Hadamard derivatives
\citep[Lemma~3.9.27]{vanderVaartWellner1996}.
The argmax functional $\phi:\mathbb{P}\mapsto
\operatorname{argmax}_{p\in[\delta,1-\delta]}q_{\mathbb{P}}(p)$
is then Hadamard differentiable at $F$ tangentially to the set of
directions $h$ for which $\phi(F+th)$ has a unique maximiser near
$p^*$ for all small $t>0$.  Under the strict local maximum condition
$q''(p^*)<0$ (Assumption A3 and Theorem~\ref{thm:asymptotics}),
this tangential Hadamard differentiability holds; the general theory
is developed in \citet{DumbgenKopp1998} and, with the weakest
possible conditions for bootstrap validity, in
\citet{FangSantos2019}.  The nonparametric bootstrap is consistent
for Hadamard differentiable functionals of the empirical process
\citep[Theorem~23.9]{vanderVaartWellner1996}: the bootstrap
distribution of $n^{2/5}(\hat p^{*,b}-\hat p^*)$ (computed under
the resampling distribution $\hat{\mathbb{P}}_n$) converges in
probability to the same Gaussian limit as $n^{2/5}(\hat p^*-p^*)$,
in the bounded Lipschitz metric on probability measures on
$\mathbb{R}$. $\square$

%% ============================================================
\begin{thebibliography}{99}

\bibitem[Azzalini and Bowman(1990)]{AzzaliniBowman1990}
Azzalini, A.\ and Bowman, A.W.\ (1990).
A look at some data on the Old Faithful geyser.
\textit{Applied Statistics}, 39, 357--365.

\bibitem[Chazal and Michel(2021)]{Chazal2018}
Chazal, F.\ and Michel, B.\ (2021).
An introduction to topological data analysis: fundamental and practical
aspects for data scientists.
\textit{Frontiers in Artificial Intelligence}, 4, 667963.

\bibitem[Chernoff(1964)]{Chernoff1964}
Chernoff, H.\ (1964).
Estimation of the mode.
\textit{Annals of the Institute of Statistical Mathematics}, 16, 31--41.

\bibitem[Cont and Tankov(2004)]{ContTankov2004}
Cont, R.\ and Tankov, P.\ (2004).
\textit{Financial Modelling with Jump Processes}.
Chapman \& Hall/CRC Press.

\bibitem[D{\"u}mbgen and Kopp(1998)]{DumbgenKopp1998}
D{\"u}mbgen, L.\ and Kopp, P.\ (1998).
Extensions of Hadamard differentiability and applications.
\textit{Statistical Inference for Stochastic Processes}, 1, 261--283.

\bibitem[Efron(1987)]{Efron1987}
Efron, B.\ (1987).
Better bootstrap confidence intervals.
\textit{Journal of the American Statistical Association}, 82, 171--185.

\bibitem[Fan and Gijbels(1996)]{FanGijbels1996}
Fan, J.\ and Gijbels, I.\ (1996).
\textit{Local Polynomial Modelling and Its Applications}.
Chapman \& Hall.

\bibitem[Fang and Santos(2019)]{FangSantos2019}
Fang, Z.\ and Santos, A.\ (2019).
Inference on directionally differentiable functions.
\textit{Review of Economic Studies}, 86, 2786--2832.

\bibitem[Good and Gaskins(1971)]{GoodSilverman1986}
Good, I.J.\ and Gaskins, R.A.\ (1971).
Nonparametric roughness penalties for probability densities.
\textit{Biometrika}, 58, 255--277.

\bibitem[Hall and York(2001)]{HallYork2001}
Hall, P.\ and York, M.\ (2001).
On the calibration of Silverman's test for multimodality.
\textit{Statistica Sinica}, 11, 515--536.

\bibitem[Hartigan and Hartigan(1985)]{HartiganHartigan1985}
Hartigan, J.A.\ and Hartigan, P.M.\ (1985).
The dip test of unimodality.
\textit{Annals of Statistics}, 13, 70--84.

\bibitem[Jiang(2017)]{Jiang2017}
Jiang, H.\ (2017).
Uniform convergence rates for kernel density estimation.
\textit{Proceedings of the 34th International Conference on Machine Learning},
PMLR 70, 1694--1703.

\bibitem[Lehmann and Casella(1998)]{LehmannCasella1998}
Lehmann, E.L.\ and Casella, G.\ (1998).
\textit{Theory of Point Estimation}, 2nd ed.
Springer.

\bibitem[M{\"u}ller and Sawitzki(1991)]{MullerSawitzki1991}
M{\"u}ller, D.W.\ and Sawitzki, G.\ (1991).
Excess mass estimates and tests for multimodality.
\textit{Journal of the American Statistical Association}, 86, 738--746.

\bibitem[Orujov(2025)]{Orujov2025EcLet}
Orujov, S.\ (2025).
Glass barriers in income distributions: nonparametric evidence.
\textit{Economics Letters}, submitted.

\bibitem[Parzen(1962)]{Parzen1962}
Parzen, E.\ (1962).
On estimation of a probability density function and mode.
\textit{Annals of Mathematical Statistics}, 33, 1065--1076.

\bibitem[Romano(1988)]{Romano1988}
Romano, J.P.\ (1988).
On weak convergence and optimality of kernel density estimates of the mode.
\textit{Annals of Statistics}, 16, 629--647.

\bibitem[Serfling(1980)]{Serfling1980}
Serfling, R.J.\ (1980).
\textit{Approximation Theorems of Mathematical Statistics}.
Wiley.

\bibitem[Silverman(1981)]{Silverman1981}
Silverman, B.W.\ (1981).
Using kernel density estimates to investigate multimodality.
\textit{Journal of the Royal Statistical Society: Series B}, 43, 97--99.

\bibitem[van der Vaart and Wellner(1996)]{vanderVaartWellner1996}
van der Vaart, A.W.\ and Wellner, J.A.\ (1996).
\textit{Weak Convergence and Empirical Processes}.
Springer.

\bibitem[Di~Martino et~al.(2025)]{ZDip2025}
Di~Martino, E., Cinelli, M., and Cerqueti, R.\ (2025).
Z-Dip: a validated generalization of the Dip Test.
\textit{arXiv:2511.01705}.

\end{thebibliography}

%% ============================================================
%% TABLES — auto-generated by src/csv_to_tex.py
%% Run `make tables` to regenerate from simulation results.
\clearpage
\begin{table}[p]
\centering
\caption{Bias and RMSE of $\hat p^*$ (percentile-point units)
across DGPs and sample sizes; $R=500$, $B=299$.}
\label{tab:bias_rmse}
\begin{tabular}{lcccccccc}
\toprule
 & \multicolumn{2}{c}{$n=250$} & \multicolumn{2}{c}{$n=500$} & \multicolumn{2}{c}{$n=1000$} & \multicolumn{2}{c}{$n=5000$}\\
\cmidrule(lr){2-3}\cmidrule(lr){4-5}\cmidrule(lr){6-7}\cmidrule(lr){8-9}
DGP & Bias & RMSE & Bias & RMSE & Bias & RMSE & Bias & RMSE\\
\midrule
A6 (unimodal, size) & --- & --- & --- & --- & --- & --- & --- & ---\\
A1 (weak, $\Delta=2.6$) & -0.005 & 0.090 & -0.004 & 0.074 & 0.001 & 0.061 & 0.002 & 0.035\\
A2 (moderate, $\Delta=2.8$) & -0.005 & 0.078 & -0.004 & 0.060 & -0.001 & 0.045 & 0.002 & 0.028\\
A3 (sharp, $\Delta=3$) & -0.006 & 0.066 & -0.004 & 0.047 & -0.000 & 0.034 & -0.000 & 0.021\\
A4 (unequal $\pi$, $\Delta=3$) & -0.015 & 0.060 & -0.017 & 0.050 & -0.011 & 0.039 & -0.009 & 0.023\\
A5 (unequal $\sigma$, $\Delta\approx2.7$) & 0.024 & 0.068 & 0.019 & 0.053 & 0.017 & 0.041 & 0.007 & 0.023\\
B1 (skew-normal) & -0.004 & 0.074 & -0.008 & 0.066 & 0.001 & 0.050 & -0.001 & 0.034\\
B2 (heavy tail $t_3$) & -0.005 & 0.075 & -0.007 & 0.056 & -0.002 & 0.049 & 0.000 & 0.027\\
B3 (Weibull mixture) & 0.027 & 0.122 & 0.036 & 0.113 & 0.047 & 0.100 & 0.063 & 0.087\\
C1 (two anti-modes) & 0.004 & 0.036 & 0.004 & 0.028 & 0.004 & 0.018 & 0.003 & 0.009\\
C2 (three anti-modes) & -0.000 & 0.031 & 0.000 & 0.021 & -0.000 & 0.015 & -0.000 & 0.006\\
C3 (bathtub) & 0.001 & 0.032 & 0.001 & 0.022 & 0.000 & 0.015 & -0.000 & 0.007\\
\bottomrule
\end{tabular}
\end{table}

\clearpage
\begin{table}[p]
\centering
\caption{Empirical convergence rates: slope of
$\log(\mathrm{RMSE})$ on $\log(n)$ (theoretical: $-0.40$).}
\label{tab:rates}
\begin{tabular}{lcc}
\toprule
DGP & Estimated slope & 95\% CI\\
\midrule
A2 & $-0.343$ & $[-0.408,\;-0.299]$\\
A3 & $-0.367$ & $[-0.484,\;-0.285]$\\
B2 & $-0.332$ & $[-0.416,\;-0.205]$\\
B3 & $-0.112$ & $[-0.166,\;-0.085]$\\
C1 (barrier 1) & $-0.481$ & $[-0.637,\;-0.346]$\\
C3 & $-0.489$ & $[-0.531,\;-0.468]$\\
\bottomrule
\end{tabular}
\end{table}

\clearpage
\begin{table}[p]
\centering
\caption{Bootstrap 95\% coverage (Cov., \%), detection rate (Det., \%),
and false detection rate (FDR, \%) under the unimodal null (DGP~A6),
across sample sizes.  Coverage is reported for DGPs with a true
anti-mode only; ``---'' denotes not applicable.
Results at $n=250$ are omitted: bootstrap coverage intervals are too
wide to be informative at that sample size (median half-width
${>}15$ percentile points).}
\label{tab:coverage}
\begin{tabular}{lccccccccc}
\toprule
 & \multicolumn{3}{c}{$n=500$} & \multicolumn{3}{c}{$n=1000$} & \multicolumn{3}{c}{$n=5000$}\\
\cmidrule(lr){2-4}\cmidrule(lr){5-7}\cmidrule(lr){8-10}
DGP & Cov. & Det. & FDR & Cov. & Det. & FDR & Cov. & Det. & FDR\\
\midrule
A6 (unimodal) & --- & --- & 13.0 & --- & --- & 9.6 & --- & --- & 1.0\\
A1 (weak) & 96.1 & 72.6 & --- & 95.2 & 72.2 & --- & 89.6 & 61.8 & ---\\
A2 (moderate) & 97.1 & 75.6 & --- & 96.0 & 73.8 & --- & 90.4 & 80.4 & ---\\
A3 (sharp) & 97.4 & 81.6 & --- & 96.6 & 83.2 & --- & 93.6 & 93.4 & ---\\
B2 (heavy tail) & 93.8 & 77.8 & --- & 94.9 & 87.0 & --- & 92.0 & 96.8 & ---\\
C1 (joint) & 93.6 & 100.0 & --- & 95.4 & 100.0 & --- & 93.0 & 100.0 & ---\\
C3 (bathtub) & 94.6 & 100.0 & --- & 95.6 & 100.0 & --- & 93.0 & 100.0 & ---\\
\bottomrule
\end{tabular}
\end{table}

\clearpage
\begin{table}[p]
\centering
\caption{Simulated rejection rates at nominal 5\% level, DGP~A2
(moderate barrier, $\Delta=2.8$); $R=500$, $B=299$.
Anti-mode test and dip test applied by the author.
Z-Dip statistic \citep{ZDip2025} applied by the author;
Di~Martino et~al.\ validate Z-Dip only on distributions with
component separation $\geq 8\sigma$ and report no results for
$\Delta=2.8\sigma$.}
\label{tab:power}
\begin{tabular}{lccc}
\toprule
$n$ & Anti-mode test & Dip test & Z-Dip\\
\midrule
 250 & 0.000 & 0.100 & 0.086\\
 500 & 0.000 & 0.216 & 0.194\\
1000 & 0.000 & 0.396 & 0.372\\
5000 & 0.688 & 0.998 & 0.998\\
\bottomrule
\end{tabular}
\end{table}


\end{document}
